\chapter{Métodos Robustos, Flexíveis e Imputação}

\section{M-estimadores robustos}

`rlm` ajusta um modelo linear com M-estimadores que reduzem o peso de
outliers. O padrão é uma perda de Huber com constante de ajuste, que
dá comportamento tipo OLS no centro e tipo L1 nas caudas.

\begin{lstlisting}[caption={Regressão robusta}]
let m_rlm = rlm(y ~ x1 + x2, df)
print(m_rlm)
\end{lstlisting}

Use `rlm` quando os dados contiverem outliers influentes, mas a maior
parte da distribuição for aproximadamente normal.

\section{Modelos Lineares Generalizados}

`glm` ajusta um preditor linear através de uma função de ligação.
Generaliza o OLS para famílias exponenciais como Poisson, Binomial e
Gamma.

\begin{lstlisting}[caption={GLM}]
let m_glm = glm(y ~ x1 + x2, df, family=poisson, link=log)
print(m_glm)
\end{lstlisting}

Famílias disponíveis incluem `gaussian`, `poisson`, `binomial`, `gamma`,
`inverse_gaussian` e `negative_binomial`. Ligações incluem `identity`,
`log`, `logit`, `probit`, `cloglog` e `inverse`.

\section{GEE}

Equações de Estimação Generalizadas ajustam um modelo marginal
considerando a correlação intra-grupo. É o análogo médio-populacional
dos modelos mistos.

\begin{lstlisting}[caption={GEE}]
let m_gee = gee(y ~ x1 + x2, df, id="grupo", family=poisson)
print(m_gee)
\end{lstlisting}

Estruturas de correlação incluem `independence`, `exchangeable`, `ar1` e
`unstructured`.

\section{Regressão Beta}

`betareg` é apropriada quando a variável dependente é limitada em
$(0,1)$.

\begin{lstlisting}[caption={Regressão Beta}]
let m_beta = betareg(share ~ x1 + x2, df)
print(m_beta)
\end{lstlisting}

Ela modela a média por uma ligação logit e a dispersão por um parâmetro
de precisão separado.

\section{GAM e lowess}

`gam` ajusta um modelo aditivo generalizado com splines de suavização
para efeitos não lineares.

\begin{lstlisting}[caption={GAM}]
let m_gam = gam(y ~ x1 + s(x2) + x3, df)
print(m_gam)
\end{lstlisting}

`lowess` é um suavizador polinomial local para exploração visual e
regressão parcial.

\begin{lstlisting}[caption={Lowess}]
let m_low = lowess(df, y, x2, frac=0.3)
print(m_low)
\end{lstlisting}

\section{Imputação múltipla}

`mice` realiza imputação múltipla por equações encadeadas. Preenche
valores ausentes usando uma sequência de modelos condicionais e retorna
um conjunto de DataFrames completos.

\begin{lstlisting}[caption={MICE}]
let m_mice = mice(df, vars=["y", "x1", "x2"], m=5)
print(m_mice)
\end{lstlisting}

As regras de Rubin combinam os resultados de cada conjunto imputado em
uma única estimativa com variância que reflete a incerteza da imputação.
